\chapter{Bent Penis (Peyronie Disease)}

\section*{Definition}
Acquired curvature of the erect penis caused by fibrous plaque formation in the tunica albuginea of the corpora cavernosa, often associated with pain and erectile dysfunction.

\section*{When to See a Doctor}

\begin{itemize}
 \item Sudden bending injury with a popping sound, immediate loss of erection, severe pain, bruising, or swelling may indicate penile fracture and requires emergency care.
 \item Arrange urologic assessment for a new or worsening curve, a palpable plaque, painful erections, penile shortening, or erectile difficulty.
\end{itemize}

\section*{How It Develops}
Peyronie disease results from repetitive microvascular trauma or buckling injury to the erect penis, triggering an abnormal wound-healing response. Fibrin deposition leads to chronic inflammation, dysregulated collagen production, and plaque formation in the tunica albuginea. The inelastic plaque prevents uniform expansion of the tunica during erection, causing curvature toward the side of the plaque.

\section*{Causes}
\begin{itemize}
 \item Peyronie disease (idiopathic, most common)
 \item Repeated minor injury or, less commonly, a major penile injury
 \item Dupuytren contracture (associated condition)
 \item Genetic predisposition (family history)
 \item Connective tissue disorders
 \item Diabetes mellitus
 \item Other fibrotic conditions and metabolic risk factors may be associated
\end{itemize}

\section*{Related Symptoms}
\begin{itemize}
 \item Erectile dysfunction
 \item Penile pain
 \item Penile shortening
 \item Dupuytren contracture
 \item Penile plaque
\end{itemize}


\section*{Medical Evaluation}
\begin{itemize}
 \item Your doctor will take a detailed history of onset, progression, pain during erections, and impact on sexual function, and perform a focused genital examination.
 \item Palpation of the penile shaft identifies the location, size, and tenderness of the fibrous plaque within the tunica albuginea.
 \item Photographs or an examination of the erect penis can help document the degree and direction of curvature. Ultrasound, sometimes with an induced erection, is used when the plaque, calcification, blood flow, or diagnosis needs further assessment.
 \item Referral to a urologist, preferably one experienced in sexual medicine, is appropriate for confirmation and treatment planning.
\end{itemize}

\section*{Treatment Options}
\begin{itemize}
 \item Penile traction therapy may modestly improve curvature or preserve length in some people when used consistently with a prescribed device; benefits and treatment duration vary.
 \item Intralesional treatments, including collagenase where available and appropriate, are specialist treatments for selected plaques and curves and require counselling about bruising, swelling, pain, and the rare risk of corporal rupture.
 \item Oral vitamin E, tamoxifen, and similar medicines have not shown reliable benefit and are not routinely recommended.
 \item Surgery is considered only when disease is stable and the deformity causes functional impairment or major distress. Options include plication, plaque incision or grafting, or penile prosthesis when erectile dysfunction does not respond to medication.
\end{itemize}

\section*{Self-Care and Home Management}
\begin{itemize}
 \item Avoid forceful bending or vigorous manipulation. Sexual activity may be continued if comfortable and safe, but stop if it causes pain or a new injury.
 \item A clinician may recommend an NSAID for pain if it is safe for you; follow the label and ask about kidney, stomach, heart, blood-thinner, and other risks.
 \item Apply gentle traction or stretching exercises only if recommended by a specialist — do not attempt manual correction of curvature.
 \item Monitor curvature by taking monthly photographs at home to objectively track progression or stability over time.
\end{itemize}

\section*{References}
\begin{thebibliography}{99}
\bibitem{bent_penis:ref1} American Urological Association. ``Peyronie\'s Disease Guideline.'' 2015. https://www.auanet.org/documents/education/clinical-guidance/Peyronies-Disease.pdf.
\bibitem{bent_penis:ref2} Ralph D, Gonzalez-Cadavid N, et al. ``The management of Peyronie\'s disease: an evidence-based review.'' \textit{Eur Urol}. 2015;68(4):656-666.
\bibitem{bent_penis:ref3} Gelbard MK, Dorey F, et al. ``Collagenase clostridium histolyticum for Peyronie\'s disease.'' \textit{N Engl J Med}. 2013;369(17):1624-1633.
\bibitem{bent_penis:ref4} Wein AJ, Kavoussi LR, et al. \textit{Campbell-Walsh Urology}, 12th ed. Elsevier, 2020.
\bibitem{bent_penis:ref5} European Association of Urology. ``Penile Curvature: EAU Guidelines on Sexual and Reproductive Health.'' 2026. https://uroweb.org/guidelines/sexual-and-reproductive-health/chapter/penile-curvature.
\end{thebibliography}
